\documentclass[11pt]{extarticle}
\usepackage[a4paper,margin=1.15cm,top=1cm,bottom=1cm]{geometry}
\usepackage{amsmath,amssymb}
\usepackage{fontspec}
\setmainfont{FreeSerif}
\setsansfont{DejaVu Sans}[Scale=0.9]
\usepackage{polyglossia}
\setdefaultlanguage{russian}
\usepackage{multicol}
\usepackage{enumitem}
\usepackage{xcolor}
\setlist{nosep,leftmargin=1.2em}
\pagestyle{empty}
\setlength{\parindent}{0pt}
\setlength{\parskip}{2pt}
\definecolor{acc}{HTML}{B4442A}
\newcommand{\blk}[2]{\par\smallskip{\sffamily\bfseries\color{acc}#1}\quad{\sffamily\small #2}\par}
\newcommand{\mol}[1]{\par\smallskip\colorbox{acc!10}{\parbox{\dimexpr\linewidth-2\fboxsep}{\textbf{#1}}}\par}
\begin{document}

{\sffamily\bfseries\Large Замена алгебры, 9 КЛ, пн 5.10}\hfill{\sffamily\small 45 мин, один план на оба урока}

\textbf{Цель Ольги:} все поняли ограниченность/неограниченность и асимптоту; строго построили отрицание (№1); в №6 \emph{сами} выдвигают гипотезы, из какого базового графика и как строить.
\textbf{Сквозная идея урока:} как функция ведёт себя \emph{на краях} — там, где $x$ огромный, и там, где знаменатель обнуляется.

\blk{0–8}{ВБРОС: игра «Ограничитель против Вылезателя»}
На доске $f(x)=\dfrac{4x}{1+x^2}$. «Я — Ограничитель: называю число $c$. Вы называете $x$. Если $|f(x)|>c$ — вы выиграли». Говорю $c=2$. Пусть пробуют: $x=1$ даёт ровно $2$ — не вылезло, $x=10$ даёт $\approx0{,}4$. Не выходит.
Теперь $f(x)=x+\dfrac4x$. Говорю $c=1000$ — отвечают $x=2000$. $c=10^6$ — снова вылезают. «У кого выигрышная стратегия?»
\textbf{Вывод на доску:} ограничена $\iff$ у Ограничителя есть выигрышная стратегия: $\exists c>0\ \forall x\in X\colon |f(x)|\le c$.
Связка со спецматом: листок «Игры», «может ли один играть так, чтобы непременно выиграть, \emph{как бы ни играл} второй» — это ровно $\exists\,\forall$.
\textbf{Задание:} «запишите формулой, что значит „выигрывает Вылезатель“. Это №1». Ответ сам не пишу.

\blk{8–20}{Решают №1–2, хожу по классу}
Ошибки в №1: меняют только знак ($\exists c\ \forall x\colon |f|>c$); пишут $\ge$ вместо $>$; теряют $X$. Подсказка: «кто теперь ходит первым?»\\
Ловушки №2: Г) ограничена сверху, но \emph{не} ограничена; Б) значения перепрыгивают интервал $(-4;4)$.

\mol{~20 мин. Молния 1 «Порядок ходов» (2 мин)}
«Пусть первым ходит Вылезатель: $\forall x\ \exists c\colon |f(x)|\le c$. Какие функции этому удовлетворяют?» — \emph{Все} ($c=|f(x)|+1$). А $\exists x\ \forall c$? — Никакие.
Мораль: порядок кванторов — это порядок ходов. Весь матанализ состоит из таких игр («для любого $\varepsilon$ найдётся $\delta$…»).

\blk{20–28}{№3}
\mol{Молния 2 «Дыра в определении» (3 мин)}
Читаю вслух опр.\,2. «По этому определению ось $Oy$ — асимптота параболы $y=x^2$? Точки графика подходят к ней сколь угодно близко!» Пусть починят: \emph{расстояние от точки графика до прямой стремится к 0, когда точка уходит по графику в бесконечность.}\\
Миф «асимптота не пересекает график» (греч. «несовпадающая»): рисую затухающую волну $y=\frac{\sin x}{x}$ — пересекает $y=0$ бесконечно много раз, а асимптота есть.\\
Обратный вопрос: $y=\sqrt x$ растёт всё медленнее. Есть горизонтальная асимптота? — Нет: функция неограничена. Связка: горизонтальная асимптота $\Rightarrow$ на этом краю функция ограничена.

\blk{28–38}{№4–5}
\mol{Молния 3 «Деление с остатком» (3 мин)}
$\frac{23}{7}=3+\frac27$. Так же: $2x+3=2(x-2)+7\ \Rightarrow\ y=2+\dfrac{7}{x-2}$. Неполное частное — асимптота $y=2$, остаток — гипербола с $k=7$, сдвинутая на 2 вправо.
Подставить $x=10^6$: $y\approx2{,}000007$. \textbf{На бесконечности побеждают старшие члены.}
Дальше: $\dfrac{x^3}{x^2+1}=x-\dfrac{x}{x^2+1}$ — наклонная асимптота $y=x$ (это вторая страница, для Цуканова).
№4 на пальцах: $\dfrac{2x+4}{x+2}$ — это константа $2$ с выколотой точкой $x=-2$ (здесь $ad-bc=0$).

\blk{38–45}{№6: гипотезы}
\textbf{Правило на доску:} модуль на $x$ (внутри) — зеркало $Oy$: правую половину копируем влево. Модуль на $y$ (снаружи) — зеркало $Ox$: нижнюю часть отражаем вверх.
Перед решением каждого пункта — вопрос: «из какого графика и в каком порядке?» Ловушка 5): $\frac{|x|+4}{x-2}$ — симметрия? Нет, знаменатель без модуля, строим по кускам.
\textbf{Финальный вопрос (2 мин, связывает урок):} «какие из шести графиков ограничены?» — только 2) и 3).

\blk{Всегда}{Подняли руку — сначала спрашиваю «из какого графика?», а не отвечаю. Кто закончил — вторая страница (не распечатана: с телефона или на доску).}

\newpage
{\sffamily\bfseries\Large Ответы}

\begin{multicols}{2}
\textbf{№1.} $\forall c>0\ \exists x\in X\colon |f(x)|>c$.

\textbf{№2.}
\begin{itemize}
\item[А)] $[-2;2]$, огр. ($1+x^2\ge2|x|$).
\item[Б)] $(-\infty;-4]\cup[4;+\infty)$, неогр. ($=x+\frac4x$).
\item[В)] $[0;3]$, огр.
\item[Г)] $(-\infty;1]$: огр. сверху, не огр.
\item[Д)] $[2;+\infty)$: огр. снизу, не огр.
\item[Е)] $(-\infty;0]\cup[12;+\infty)$, неогр. ($=x+6+\frac9x$).
\end{itemize}

\textbf{№3.} $x=0$ и $y=0$: при $|x|\to\infty$ $\frac kx\to0$; при $x\to0$ $\left|\frac kx\right|\to\infty$.

\textbf{№4.} $c=0$ — линейная функция (прямая). $ad-bc=0$ — числитель пропорционален знаменателю, $y\equiv\frac ac$ с выколотой точкой $x=-\frac dc$.

\textbf{№5.} $y=2+\frac{7}{x-2}$, $k=7$. Асимптоты $x=2$, $y=2$; центр симметрии $(2;2)$. $D\colon x\ne2$, $E\colon y\ne2$. Убывает на $(-\infty;2)$ и на $(2;+\infty)$ — \emph{не} на всей области! Нуль $x=-1{,}5$, $y(0)=-1{,}5$.

\textbf{№6} (цепочки):
\begin{enumerate}[label=\arabic*)]
\item $\frac4{|x+2|}=\left|\frac4{x+2}\right|$: $\frac4x\to$ влево на 2 $\to$ низ вверх. Неогр.
\item $\frac4{|x|+2}$: $\frac4{x+2}$, правую половину отразить влево. Огр., $E=(0;2]$.
\item $2+\frac1{|x|+2}$: $\frac1x\to$ влево 2 $\to$ вверх 2 $\to$ правую отразить. Огр., $E=(2;2{,}5]$.
\item $\left|2+\frac1{x+2}\right|$: $\frac1x\to$ влево 2 $\to$ вверх 2 $\to$ низ вверх. Нуль $-2{,}5$. Неогр.
\item Не симметричен! $x\ge0$: $1+\frac6{x-2}$; $x<0$: $-1+\frac2{x-2}$. Неогр.
\item $\left|2+\frac3{|x|-1}\right|$: $\frac3x\to$ вправо 1 $\to$ вверх 2 $\to$ правую отразить $\to$ низ вверх. Неогр.
\end{enumerate}

\columnbreak
{\sffamily\bfseries Вторая страница (Цуканов)}

\textbf{109} (верт.; гориз./накл.):
а) $x=\pm1$; $y=0$.\quad б) $x=0,-1,-2$; $y=0$.\\
в) $x=\pm\sqrt2,\pm2$; $y=0$.\quad г) $x=0$; $y=x$.\\
д) нет; $y=x+1$ ($=x+1-\frac1{x^2+1}$).\\
е) $x=-2$; $y=x-4$ ($+\frac{11}{x+2}$).\\
ж) $x=0$; $y=2x+1$ ($+\frac1{x^3}$).\\
з) $x=-1$; $y=x-2$ ($+\frac{3x+2}{(x+1)^2}$).\\
и) нет; $y=x$ ($-\frac{x}{x^4+1}$).

\textbf{110} (знак остатка):
а) $y=x$, остаток $\frac1{x^2}>0$: сверху с обеих сторон.\\
б) $y=x-1$, остаток $\frac{3x+5}{(x+2)^2}\approx\frac3x$: при $+\infty$ сверху, при $-\infty$ снизу.\\
в) при $+\infty$: $y=x-1$, сверху; при $-\infty$: $y=-x+1$, сверху.

\textbf{112} (асимптоты; корни):
а) $x=0$, $y=x$; $\pm1$.\quad б) $x=-1$, $y=x-1$; нет.\\
в) \emph{ловушка:} $=x+1$ при $x\ne1$ — асимптот нет, прямая с дыркой; корень $-1$.\\
г) $x=1$, $y=x+2$; $0,-1$.\quad д) $x=1$, $y=-1$; $0$.\\
е) $y=x$; $0$.\quad ж) $x=\pm1$, $y=x$; $0$.\\
з) $x=\pm4$, $y=-4x$; $0$.\quad и) $x=\pm4$, $y=-4x$; $0,\pm\frac32$.\\
к) $x=\pm5$, $y=x$; $0,\pm4$.

\medskip
{\sffamily\bfseries Метод для всех: } поделить с остатком. Неполное частное — асимптота; знак остатка — сверху или снизу; корни знаменателя — вертикальные асимптоты (если не сократились).
\end{multicols}
\end{document}
